\documentclass[12pt]{article}
\usepackage[margin=0.75in]{geometry}
\usepackage{fontspec}
\setmainfont{Latin Modern Roman}
\usepackage{microtype}
\usepackage{multicol}
\usepackage{fancyhdr}
\usepackage{titlesec}
\usepackage{enumitem}
\usepackage{xcolor}
\usepackage[hidelinks]{hyperref}
\setlength{\columnsep}{0.28in}
\setlength{\parindent}{1.1em}
\setlength{\parskip}{0.32em}
\setlength{\headheight}{15pt}
\titleformat{\section}{\large\bfseries\scshape}{}{0pt}{}
\titleformat{\subsection}{\normalsize\bfseries}{}{0pt}{}
\pagestyle{fancy}
\fancyhf{}
\fancyfoot[C]{\thepage}
\renewcommand{\headrulewidth}{0.4pt}
\newcommand{\focusbox}[1]{\par\vspace{0.4em}\noindent\begingroup\setlength{\fboxsep}{6pt}\colorbox{gray!12}{\begin{minipage}{\dimexpr\linewidth-2\fboxsep\relax}#1\end{minipage}}\endgroup\par\vspace{0.5em}}

\fancyhead[L]{Whately: Disciplined Thinking}
\fancyhead[R]{Student Reading}
\begin{document}
\begin{center}
{\LARGE\bfseries What It Means to Think Straight}\par\vspace{0.25em}
{\large\scshape Week 19: Final Portfolio Defense}\par\vspace{0.5em}
{Defending intellectual growth with Whately and portfolio evidence}\par\vspace{0.6em}
\rule{0.82\textwidth}{0.4pt}
\end{center}

\section*{Central Question}
What does it mean to be a disciplined thinker?

\section*{Vocabulary Preview}
\begin{multicols}{2}
\begin{itemize}[leftmargin=*, itemsep=0.18em]
\item \textbf{Defense}: an argument that supports a judgment with reasons and evidence.
\item \textbf{Controlling claim}: the precise thesis the defense will establish.
\item \textbf{Portfolio evidence}: an artifact used to support a claim about growth.
\item \textbf{Commentary}: explanation of how evidence supports the claim.
\item \textbf{Warrant}: the principle connecting evidence to a conclusion.
\item \textbf{Qualification}: a limit that makes a claim more accurate rather than weaker.
\item \textbf{Counterpressure}: evidence or an objection that tests the thesis.
\item \textbf{Transfer claim}: an account of how a habit applies beyond this course.
\item \textbf{Disciplined thinker}: one who repeatedly clarifies, tests, and judges rather than merely sounding logical.
\item \textbf{Intellectual humility}: willingness to revise when reasons or evidence require it.
\end{itemize}
\end{multicols}

\section*{Introduction}
The final defense is not a victory speech and not a list of completed assignments. It is an argument. You will define disciplined thinking, use at least three artifacts as evidence, explain what each artifact demonstrates, examine a limit or counterpressure, and show how one habit transfers beyond the unit. At least one artifact must be revised, and at least one must involve \textit{Macbeth} or \textit{Julius Caesar}.

Whately's course ends in modest power. Logic does not discover every fact, guarantee wise motives, or replace moral judgment. It does train habits: define terms, state claims clearly, test whether conclusions follow, inspect evidence and sources, expose hidden assumptions, distinguish rhetoric from proof, and identify what remains genuinely disputed.

\focusbox{\textbf{Source note:} This final reading synthesizes the course's adapted Whately sequence. The defense architecture is an original course scaffold. The aim is a warranted, qualified claim about growth---not a claim of perfect rationality.}

\begin{multicols}{2}
\section*{Reading}

\subsection*{1. Define Before You Defend}
``A disciplined thinker thinks carefully'' is circular and too vague. Define the phrase with observable habits. A strong definition might include stable terms, clear claims, tested inference, proportionate evidence, fair attention to objections, and revision when the case fails.

Your definition should guide the whole defense. If intellectual humility matters in the thesis, at least one artifact should show an actual revision, qualification, or corrected judgment.

\subsection*{2. Make a Controlling Claim}
A thesis should be arguable and specific. It may name two or three habits and a pattern of growth. Avoid claiming that the course made you immune to error. A defensible thesis explains what changed and how the portfolio makes that change visible.

The thesis is not ``I completed many assignments.'' Completion is a fact about work submitted. The defense concerns the quality of judgment recorded in that work.

\subsection*{3. Use Artifacts as Evidence}
For each artifact, provide enough detail that a reader can verify the claim:

\begin{itemize}[leftmargin=*]
\item identify the assignment or passage;
\item quote or accurately describe the original reasoning;
\item name the logical habit at work;
\item explain the before-and-after change when revised;
\item connect the artifact to the thesis.
\end{itemize}

An artifact does not speak for itself. Commentary must show why a narrower conclusion, clarified term, stronger source test, or repaired inference counts as evidence of disciplined thinking.

\subsection*{4. State the Warrant}
The warrant is the bridge between artifact and thesis. If an artifact shows that you changed ``all students'' to ``students in this survey,'' the warrant might be: disciplined thinkers size conclusions to the evidence available. Without the bridge, the defense becomes a scrapbook with praise attached.

\subsection*{5. Test Your Own Defense}
Apply logic to the final argument itself. Ask:

\begin{itemize}[leftmargin=*]
\item Are my key terms stable?
\item Do three selected artifacts represent a pattern, or only isolated successes?
\item Have I proved growth or merely effort?
\item Is my Shakespeare example analyzed rather than summarized?
\item What counterexample, remaining weakness, or qualification should I acknowledge?
\end{itemize}

A qualification can strengthen the thesis. ``I now always reason perfectly'' is incredible. ``I have learned to slow down at three recurring pressure points, though source judgment remains difficult'' is testable and honest.

\subsection*{6. Use Shakespeare as a Judgment Mirror}
The plays give negative examples of disciplined thought. Macbeth often mistakes equivocal security for certainty; Brutus can turn possibility into political necessity. Do not claim that you are morally superior to tragic characters. Compare habits of judgment: what warning did the character ignore, and what course habit helps a reader notice the failure?

The literary artifact must do argumentative work. It may show why term clarification, source testing, evidence proportion, or hidden-assumption analysis matters.

\subsection*{7. Make a Transfer Claim}
Name a setting beyond this unit: research, workplace decisions, online claims, voting, family conflict, theology, literature, or your own writing. Then describe a concrete action. ``Logic helps in life'' is too general. ``Before sharing a dramatic statistic, I will check the sample and whether the conclusion exceeds it'' is a transfer claim.

\subsection*{8. The Six-Part Defense}
\begin{enumerate}[leftmargin=*,itemsep=0.12em]
\item \textbf{Define} disciplined thinking in observable habits.
\item \textbf{Claim} a specific pattern of intellectual growth.
\item \textbf{Present} at least three portfolio artifacts.
\item \textbf{Explain} the warrant connecting each artifact to the claim.
\item \textbf{Test} the defense with a counterpressure and qualification.
\item \textbf{Transfer} one habit to a concrete setting beyond the course.
\end{enumerate}

\section*{Final Defense Requirements}
\begin{itemize}[leftmargin=*]
\item a clear thesis about disciplined thinking;
\item accurate use of at least three logic concepts;
\item evidence from at least three portfolio artifacts;
\item at least one revised artifact;
\item at least one \textit{Macbeth} or \textit{Julius Caesar} artifact;
\item direct reflection on growth, one qualification, and one concrete transfer.
\end{itemize}

\section*{Reading Focus}
\begin{enumerate}[leftmargin=*]
\item Why is the final defense an argument rather than a celebration?
\item What makes a definition of disciplined thinking observable?
\item What is the difference between portfolio evidence and commentary?
\item Define warrant and give a one-sentence example.
\item Why can a qualification strengthen a defense?
\item What argumentative work should the Shakespeare artifact do?
\item Copy the six-part defense.
\item What makes a transfer claim concrete?
\end{enumerate}

\section*{Handout Summary}
Write 5--7 sentences explaining how to build and test a portfolio defense. Use the words \textit{controlling claim}, \textit{artifact}, \textit{warrant}, \textit{qualification}, and \textit{transfer}.

\begingroup\small
\section*{Defense Preparation}
Before drafting, place the three selected artifacts beside this page. Use the prompts below to test whether the defense is ready:
\begin{itemize}[leftmargin=*,itemsep=0.05em,topsep=0.1em]
\item Define disciplined thinking through observable habits.
\item Name the feature in each artifact that supports the controlling claim and state its warrant.
\item Identify the revised artifact and the Shakespeare judgment mirror.
\item Write the strongest fair counterpressure and the qualification it requires.
\item Name the setting, reasoning pressure, and action in the transfer claim.
\end{itemize}
\endgroup
\end{multicols}
\end{document}
